\documentclass[11pt]{beamer}

\usepackage[utf8]{inputenc}
\usepackage[T1]{fontenc}
\usepackage{lmodern}
\usepackage[french]{babel}
\usetheme{Copenhagen}
%\usecolortheme{beaver}
\useoutertheme{split} % ???
\usefonttheme{serif}

\setbeamertemplate{theorems}[numbered] % to number
\setbeamertemplate{footline}[frame number]

% Formules mathématiques
\usepackage{amsmath}
\usepackage{amsfonts}
\usepackage{amssymb}
\usepackage{amsthm}
\usepackage{nicefrac}
\usepackage{mathrsfs}

\usepackage{tikz}
\usetikzlibrary{arrows.meta}
\usepackage[linguistics]{forest}
\usepackage{graphicx}
\graphicspath{ {../simulation/figures/}}
\usepackage{xcolor}


\author{Hugo \textit{Vangi}lluwen Lilian Etchaïde}
\title{Indice de pouvoir dans l'Union Européenne}
\subtitle{TIPE}
\date{\today}
\subject{Systèmes de vote}

\begin{document}
	%\logo{}
	%\institute{}
	%\setbeamercovered{transparent}
	%\setbeamertemplate{navigation symbols}{}
	
	\begin{frame}[plain]
		\maketitle
	\end{frame}
	
	\begin{frame}
		\small \tableofcontents % [hideallsubsections]
	\end{frame}
	
	\AtBeginSection{
		\begin{frame}
			\small \tableofcontents[currentsection]
		\end{frame}
	}
	
	\section{Défauts de l'Union Européenne}
	\begin{frame}{Titre}
		Lilian
	\end{frame}
	
	\section{Théorie du choix social ou théorie des jeux ???}
	
	\subsection{Jeu coopératif}
	\begin{frame}
		\begin{definition}
			\textbf{Jeu de vote pondéré} : Pour / Contre \\
			\begin{equation*}
				(q, \left(w_i\right)_{i \in \mathcal{V}}) % quota, poids
			\end{equation*}
			$\mathcal{V}$ ensemble des votants
		\end{definition}
		\uncover<2>{\begin{definition}
				\textbf{Coalition gagnante} $S$
				\begin{equation*}
					\sum_{i \in S} w_i \geqslant q
				\end{equation*}
				% sinon perdante
		\end{definition}}
	\end{frame}
	
	\begin{frame}
		\begin{equation*}
			v(S) = \left\{ \begin{array}{cl}
				1 &\text{ si } \sum_{i \in S} w_i \geqslant q \\
				0 &\text{ sinon}
			\end{array} \right.
		\end{equation*}
		\uncover<2>{\begin{definition}
				Un électeur \textbf{décisif} est nécessaire à une coalition gagnante. \\
				Décidabilité $d_i(v)$ pour $i \in \mathcal{V}$ :
				\begin{equation*}
					d_i(v) = \sum_{S \subset \mathcal{V}, \ i \in S} v(S) - v(S \setminus \{i\})
				\end{equation*}
		\end{definition}}
	\end{frame}
	
	\subsection{Indice de Banzhaf}
	\begin{frame}
		\begin{definition}
			Indice de Banzhaf normalisé :
			\begin{equation*}
				B_i(v) = \frac{d_i(v)}{\sum_{j=1}^{n} d_j(v)}
			\end{equation*}
			Indice de Banzhaf relatif :
			\begin{equation*}
				B_i(v) = \frac{d_i(v)}{2^{|\mathcal{V}| - 1}}
			\end{equation*}
		\end{definition}
		\uncover<2>{L'algorithme naïf a une complexité de $O(n 2^n)$. ???}
	\end{frame}
	
	\begin{frame}
		\underline{Algorithme de Brams-Affuso} \\
		\uncover<2-5>{
			$d_i$ : décidabilité du votant $v$ \\
			$d_k^i$ : nombre de coalitions n'incluant pas $i$ et de poids total $k$
		}
		\uncover<3-5>{\begin{equation*}
				d_i = \sum_{k=q-w_i}^{q-1} d_k^i
				\qquad \qquad
				W = \sum_{i \in \mathcal{V}} w_i
		\end{equation*}}
		\uncover<4-5>{\begin{equation*}
				\prod_{j \in \mathcal{V}, j \neq i} \left(1 + X^{w_j}\right)
				= \sum_{k=0}^{W-w_i} d_k^i X^k
		\end{equation*}}
		\uncover<5>{Complexité en $O(n^3)$}
	\end{frame}
\end{document}